\section{TOTEモデルと未解決問題のインストール：$\aleph$ によるループの無限回転}

TOTEモデル（Test-Operate-Test-Exit）とは、1960年にジョージ・ミラー（George A. Miller）、ユージン・ギャランター（Eugene Galanter）、カール・プリブラム（Karl H. Pribram）の3名が著書『プランと行動の構造（Plans and the Structure of Behavior）』において提唱した、生体の行動や問題解決プロセスを記述する目的論的フィードバック制御モデルである。

ここで、生体の観測・計算能力を超え、その発生理由や因果構造を確定・了解することができない現象kjを\textbf{「Unlimited Unknown（超限の不可知）」}と定義し、記号 $\aleph$（アレフ）を割り当てる。

生体がこの不可知現象 $\aleph$ を知覚したとき、神経系疑似ベイズ推論機はいかなる既成のナラティブ仮説 $\hat{\Omega}$ にも収束できず、Exit条件（終了判定）を満たせない開ループのTOTEプロセスへと突入する。
すなわち、確信 $\Omega$ の確定が強制保留され、全神経系が解を求めて探索を続ける変性意識状態 $S(\mathrm{ALTER})$ が不可避に立ち上がる。
\begin{equation}
    \aleph :\xrightarrow[]{\hat{\Omega}} \mid \quad \implies \quad S(\mathrm{ALTER})
\end{equation}

このとき、推論ループが未完了のまま別の突発的な外的刺激によって処理が強制中断（インターラプト）された場合、その課題は「未完了の開放ループ（Open Loop）」として神経系のレジスタに保持・インストールされる。
その結果、後日それに類似した微小な入力ノイズやきっかけ（エビデンス $e$）が生体に降り注いだ瞬間、凍結されていた探索プロセスが即座に再点火され、生体は再び $S(\mathrm{ALTER})$ へと自動的に引き戻されるのである。


